\begin{pastexam}[label={px:df3-101-ds-01}]{地方三等資訊處理 101 年 第一題}
（申論；20 分）

在某些複雜的應用上（例如人工智慧（artificial intelligence）的應用）經常會使用經驗法則搜尋（heuristic
search），在搜尋空間的儲存上可選擇樹（tree）或圖（graph），這些樹或圖都是由點（node）與邊（edge）所構成的。
而經驗法則搜尋無非就是從起始狀態（initial node）透過經驗法則（heuristic）的指引能順利找到目標狀態（goal
node）。
\begin{itemize}[leftmargin=2.2em]
\item[㈠] 請問樹與圖最大的差異點為何？（5 分）
\item[㈡] 請問搜尋空間用樹來表示，其優缺點為何？（5 分）
\item[㈢] 請問搜尋空間用圖來表示，其優缺點為何？如果一定要用圖來表示搜尋空間，如何避免所帶來的缺點？
（10 分）
\end{itemize}
\end{pastexam}

\begin{pastsoln}
答案要點：㈠ 樹沒有環、任兩點之間只有一條路徑（每個點只有一個父節點），圖可以有環、兩點之間可以有多條路徑；
㈡ 樹不會繞圈、實作簡單，但同一個狀態會重複出現，空間可能大很多；㈢ 圖每個狀態只存一次、比較省，但會繞圈
和重複展開，要用「已拜訪集合」（closed list）記住看過的狀態來避免。

\textbf{㈠} 樹是\textbf{連通且沒有環}的圖（資結書第~7.1~節、第~21.1~節）。 等價的說法：$n$ 個點的樹恰有
$n - 1$ 條邊；任兩點之間\textbf{恰有一條}簡單路徑；有根樹裡根以外的每個點\textbf{只有一個父節點}。 一般的圖
沒有這些限制：可以有環、兩點之間可以有好幾條路徑、一個點可以由好幾個點連進來（在搜尋裡就是一個狀態可以由
好幾個不同的狀態到達），也不一定連通。 所以最大的差異是\textbf{有沒有環、路徑是否唯一}。

\textbf{㈡ 用樹表示搜尋空間}：每個節點代表「從起始狀態走到這裡的一條路徑」，同一個狀態從不同路徑到達時
是不同的節點。
\begin{itemize}
\item 優點：沒有環，搜尋（深度優先、廣度優先、最佳優先）只要一直往子節點展開，不會繞回原來的地方，不必
檢查重複；每個節點只有一個父節點，找到目標後沿父指標往回走就得到解的路徑；實作和分析都簡單。
\item 缺點：同一個狀態會在樹裡出現很多次，重複展開它底下相同的子樹，浪費時間和記憶體；狀態圖有環時，
展開成的樹是無窮的（例如可以來回走的兩個狀態），深度優先搜尋可能一直往下走不回來。 例如在方格上每步可以往上、下、左、
右走，$d$ 步內到得了的格子只有 $O(d^2)$ 個，展開成樹卻有 $4^d$ 個節點。
\end{itemize}

\textbf{㈢ 用圖表示搜尋空間}：每個節點代表一個狀態，到達它的不同路徑共用這個節點。
\begin{itemize}
\item 優點：每個狀態只存一次，空間小；不同路徑共用同一個狀態底下的部分，不必重複搜尋；能表示可以來回或
循環的問題。
\item 缺點：有環，天真地照樹的方式展開會繞圈，甚至無窮迴圈；同一個狀態會被不同路徑反覆走到，重複展開；
同一個狀態有好幾個前驅，要另外記住「目前最好的那條路從哪裡來」，路徑才還原得回來。
\end{itemize}
\textbf{避免的方法}（圖搜尋，graph search）：
\begin{enumerate}[label=(\arabic*)]
\item 多維護一個\textbf{已拜訪集合}（closed list，通常用雜湊表，查詢 $O(1)$），存已經展開過的狀態；
產生新狀態時先查，已經展開過或已經在待展開串列（open list）裡的就不再加入。 這和 BFS、DFS 替頂點著色、
只拜訪白色頂點是同一個想法（資結書第~21.3、21.4~節），每個狀態最多展開一次，不會繞圈。
\item 每個狀態記下目前找到的最小成本 $g$ 和父指標。 經驗法則搜尋（如 A*）遇到已經看過的狀態時，若新路徑
的成本更小，就更新 $g$ 和父指標（必要時把它重新放回 open list），保證找到的仍是最好的路徑；經驗函數滿足
一致性（consistent）時，第一次展開某狀態時的成本就是最小的，可以直接略過重複的狀態。
\item 深度優先搜尋若不想存整個已拜訪集合，至少檢查「目前這條路徑上」有沒有出現過這個狀態，避免走回
自己的祖先形成環，或設定深度上限。
\end{enumerate}
\end{pastsoln}
